\documentclass{article}
\usepackage[left=3cm,top=3cm,right=2cm,bottom=2cm]{geometry}
\usepackage{amssymb}
\usepackage{amsmath}
\usepackage{cancel}
\usepackage[hidelinks]{hyperref}
\usepackage{tikz}
\usepackage{booktabs}
\usepackage{float}
\usepackage{graphicx}
\usepackage[brazilian]{babel}

\graphicspath{ {../images/} }


\title{(Universidade de São Paulo)\\
Problemas de Valor Inicial}

\author{Métodos Numéricos em Equações Diferenciais\\
Docente: Fabrício Simeoni de Sousa\\[.2cm]
\begin{tabular}{r@{\ }c@{\ }l}
    Kainã Alves Tureso      &-- &15466391\\
    Vinícius de Sá Ferreira &-- &15491650
\end{tabular}}

\date{\today}

\begin{document}
    \maketitle

    \tableofcontents

    \section{Resumo}

    \section{Problema}
    No presente estudo, buscamos por resolver dois Problemas de Valor Inicial (PVI) com métodos numéricos. O primeiro problema consiste na solução da equação que rege o movimento de um pêndulo simples sem forças dissipativas, dada por
    \begin{align}
        \begin{cases}
            q''(t) + \sin(q(t)) &= 0 \\
            q(0) = \frac{\pi}{4} & q'(0) = 0
        \end{cases} \label{eq:pendulo_simples}
    \end{align}

    em que $q$ é o ângulo entre a haste e a vertical e $t \in [0,30]$

    O segundo problema é o pêndulo duplo sem forças dissipativas, descrito pelas equações
    \begin{align}
        \begin{cases}
            q_1'(t) = \dfrac{p_1 - p_2 \cos(q_1 - q_2)}{2 - \cos^2(q_1 - q_2)} \\
            q_2'(t) = \dfrac{2p_2 - p_1 \cos(q_1 - q_2)}{2 - \cos^2(q_1 - q_2)} \\
            p_1'(t) = -A_1 + A_2 - 2\sin(q_1)\\
            p_2'(t) = A_1 - A_2 - \sin(q_2)
        \end{cases} \label{eq:pendulo_duplo}
    \end{align}
    em que $q_1$ e $q_2$ são os ângulos das hastes do pêndulo conectado ao ponto fixo e o pêndulo conectado à primeira haste com a vertical, respectivamente e $p_1$ e $p_2$ são os momentos dos pêndulos. Além disso, temos
    \[
        A_1 = \dfrac{p_1p_2\sin(q_1 - q_2)}{1 + \sin^2(q_1 - q_2)} 
    \]
    \[
        A_2 = \dfrac{[p_1^2 + 2p_2^2 - 2p_1p_2 \cos(q_1 - q_2)] \sin(q_1 - q_2) \cos(q_1 - q_2)}{(1+ \sin^2(q_1 - q_2))^2}
    \]

    Os valores iniciais a serem utilizados são $q_{1; 0}$, $q_{2;0}$ $\in (0,\pi)$ e $p_{1;0} = p_{2;0} = 0$. Por fim, $t\in[0,60]$

    \section{Discretização}
    Transformaremos os PVIs \ref{eq:pendulo_simples} e \ref{eq:pendulo_duplo} para a forma vetorizada
    \begin{align}
        \mathbf{u}'(t) = \mathbf{f}(\mathbf{u}(t), t) \label{eq:forma_padrao}
    \end{align}
    em que $\mathbf{u} = (u_0(t), \dots, u_m(t))$ e $\mathbf{f}(\mathbf{u}(t), t) = (f_0(\mathbf{u}(t),t), \dots, f_m(\mathbf{u}(t),t))$.
    
    \subsection{Pêndulo simples}
    Aplicaremos os métodos de Euler explícito, Euler implícito e Runge-Kutta padrão de 4 estágios para discretizar o problema. Primeiro, reescrevemos o PVI \ref{eq:pendulo_simples} para o formato em \ref{eq:forma_padrao}:
    \begin{align}
        \mathbf{u}(t) = \begin{bmatrix}
            q(t) \\
            q'(t)
        \end{bmatrix}; &&
        \mathbf{f}(\mathbf{u}(t),t) = \begin{bmatrix}
                q'(t) \\
                -\sin(q(t))
        \end{bmatrix} \notag
    \end{align}

\end{document}